\section{Primitive coordinates and uniform recovery}\label{sec:one-slope-all-indices}
This section compares the reduced supports of scaled torsion coordinates, proves the uniform bound and one-slope results of Section~1, and isolates the remaining central comparison at wild indices.

Throughout, $p\ge5$. After completing the maximal unramified base, take the common tame good field $L/K$ of degree $m=p^2-1$. Write $R=\Ocal_L$, choose $\tau(\pi)=\zeta\pi$ with $\zeta$ of order $m$, and normalize $v_L(\pi)=1$. Put $e=e_K$ and $v=v_p(e)$, so $v_L(p)=me$. We use the good formal models and exact tame eigenparameters of Sections~\ref{sec:quotients} and~\ref{sec:restricted}. There is no assumption that $p\nmid e$.

\subsection{Minimum values and the scaled graph}
The torsion valuations determine the multiplicities of zero in the special coordinate algebras.

\begin{lemma}\label{os-lem:minima}
For a one-slope model, every primitive $p^n$-point has value $e/p^{2n-2}$; exactly $p^{2n-2}$ points of the full $p^n$-torsion, including zero, have larger value. For a two-slope model of parameter $z$, the minimum nonzero value is $z/p^{2n-1}$; exactly $p^{2n-1}$ points, including zero, have larger value.
\end{lemma}
\begin{proof}
In the one-slope case the first-level values are all $e$. For a nonzero torsion target of value at most $e$, the degree-$p$ and degree-one coefficients of its division equation lie above the segment ending in degree $p^2$. All $p^2$ divisions therefore have the target value divided by $p^2$. Induction gives the primitive values, and the higher-valued points are exactly the $p^{n-1}$-torsion.

In the two-slope case the first-level minimum is $z/p$, attained at $p^2-p$ points. Every minimum target divides on the degree-$p^2$ branch. A larger target below the middle Newton break has strictly larger division values, and above that break all division values are at least $z/p$. Divisions of zero cannot lower the minimum. Thus the minimum occurs at $(p^2-p)p^{2n-2}$ points, leaving $p^{2n-1}$ larger-valued points. The Newton range of Lemma~\ref{rc-lem:newton} is preserved because multiplication by $p$ strictly raises a positive nonzero torsion value. These formulas agree with the elliptic valuation and counting formulas of \citet[Theorem~4.4 and Corollaries~4.5--4.6, pp.~10--11]{Smith2023}; the argument here applies to the formal models under consideration.
\end{proof}


Let $f:G[p^n]\simeq H[p^n]$ and put $D=p^{2n-1}$. Make the Eisenstein extension
\[
 R'=R[\varpi]/(\varpi^D-\pi).
\]
For a one-slope coordinate assign exponent $a=pe$, and for a two-slope coordinate assign $a=z$. If $x,y$ are the coordinates in the generic graph algebra, define
\[
 X=x/\varpi^a,\qquad Y=y/\varpi^b,\qquad B'=R'[X,Y].
\]
All geometric values are integral by Lemma~\ref{os-lem:minima}. Their characteristic polynomials have coefficients in the integrally closed DVR $R'$, so $B'$ is finite and torsion-free, hence finite free. Scaling is invertible generically; its rank is therefore $p^{2n}$.

Every nonlinear term of total degree $r$ in the scaled first formal law has the extra factor $\varpi^{a(r-1)}$, and similarly for the second law. The series converge in the finite complete algebra $B'\otimes_{R'}B'$. They give integral coproducts, inverses and counits. Thus
\[
 B_0=B'/\varpi B'
\]
is a Hopf algebra generated by primitive $X,Y$. The classical primitive-presentation lemma, proved in full as Lemma~\ref{pg-lem:primitives}, gives
\begin{equation}\label{os-eq:primitive-dimension}
 V=\operatorname{span}_k\{X,X^p,\ldots,Y,Y^p,\ldots\},
 \qquad \dim_k V=2n.
\end{equation}

The group scheme $\mu_D$ acts by $\varpi\mapsto t\varpi$, fixing the original graph coordinates. It preserves $B'$ and gives a $\Z/D\Z$ grading on $B_0$, including in characteristic $p$. The weights of the Frobenius powers are
\[
 \operatorname{wt}(X^{p^j})=-ap^j,\qquad
 \operatorname{wt}(Y^{p^j})=-bp^j\pmod D.
\]
For an exponent of $p$-adic valuation $s$, there are exactly $\max(0,2n-1-s)$ distinct nonzero weights before the chain reaches zero.

\subsection{Reduced supports and stable Frobenius spans}
We show that passage to the graph preserves every reduced coordinate value; the resulting support sizes distinguish the two polygon types.

\begin{lemma}\label{os-lem:stable-spans}
In the scaled graph, the stable cyclic Frobenius span of a one-slope coordinate has dimension two; that of a two-slope coordinate has dimension one. Every Frobenius power of either coordinate is nonzero.
\end{lemma}
\begin{proof}
For one coordinate alone, $C'=R'[X]$ is monogenic of rank $p^{2n}$. Its generic coordinate has that many distinct values. The monic polynomial with those roots has integral coefficients; monic division proves that it is the exact presentation of $C'$. The scaled group law makes the special coordinate primitive.

Its monic special polynomial $P(U)$ is additive. Indeed $P(U+W)-P(U)-P(W)$ vanishes in the tensor square of the quotient and has degree below $\deg P=p^{2n}$ in each variable. The normal monomial basis makes it identically zero. The binomial-coefficient criterion then leaves only $p$-power exponents.

For a one-slope coordinate, Lemma~\ref{os-lem:minima} makes the zero-root multiplicity exactly $p^{2n-2}$. Hence
\[
 P(U)=U^{p^{2n}}+a_1U^{p^{2n-1}}+a_0U^{p^{2n-2}},
 \qquad a_0\ne0.
\]
Over the perfect field $k$, this is a $p^{2n-2}$th power of a separable additive polynomial of degree $p^2$. Its distinct roots form a two-dimensional $\F_p$-space. For a two-slope coordinate the zero multiplicity is $p^{2n-1}$, so
\[
 P(U)=U^{p^{2n}}+b_0U^{p^{2n-1}},\qquad b_0\ne0,
\]
and its distinct roots form an $\F_p$-line.

The extension $C'\subset B'$ is finite integral. Lying over makes its map on spectra surjective, including on the closed fibre. Thus all these distinct coordinate values survive in the graph, even though the map of special coordinate algebras need not be injective.

Let $F$ be $p$th power on $V$. On the two-dimensional reduced coordinate support, the functions $U,U^p$ are linearly independent: a relation $U^p=cU$ could have at most $p$ roots. The separable degree-$p^2$ relation expresses $U^{p^2}$ in their span with nonzero linear coefficient, making that span $F$-bijective. Some Frobenius power kills the nilradical of $B_0$. Applying it to the support relation gives a two-dimensional stable cyclic span in $B_0$ itself; evaluation proves that its dimension cannot drop. On the one-dimensional coordinate support the same argument gives stable dimension one. Nonzero support values show that every Frobenius power is nonzero. The stable dimensions remain the same for every tail of a chain.
\end{proof}

\subsection{Comparing different polygons}
The different stable dimensions prevent a common nonzero weight from identifying the two coordinate chains.

\begin{proof}[Proof of Theorem~\ref{thm:mixed-polygon-obstruction}]
Let $X$ be the one-slope coordinate and $Y$ the two-slope coordinate, with $w=v_p(z)$. The respective numbers of distinct nonzero weights are
\[
 d_X=\max(0,2n-2-v),\qquad
 d_Y=\max(0,2n-1-w).
\]
At an overlapping nonzero weight the two vectors are independent. Otherwise $X^{p^i}=cY^{p^j}$ with $c\ne0$, and successive Frobenius powers would identify stable cyclic spans of dimensions two and one, contrary to Lemma~\ref{os-lem:stable-spans}. Every remaining nonzero weight contributes a nonzero vector, and different weights are independent. Weight zero contains the two-dimensional stable span of $X$. Consequently
\begin{equation}\label{os-eq:dimension-inequality}
 2n\ge d_X+d_Y+2.
\end{equation}
If $w=0$, then $d_Y=2n-1$, an impossibility. Thus $w\ge1$. If $v\le2n-2$ and $w<2n-1$, substitution gives $v+w\ge2n-1$. If $w\ge2n-1$, that inequality is automatic. When $v>2n-2$, the condition $w\ge2n-1-v$ is already weaker than $w\ge1$. Therefore
\[
 w\ge\max\{1,2n-1-v\}.
\]
Finally $z\ge p^w$ and $z<pe$ give
\[
 e>z/p\ge p^{w-1}\ge p^{2n-2-v}.
\]
This includes negative exponents on the right and proves both assertions without a hypothesis on the special characters.
\end{proof}

\subsection{Comparing two one-slope models}
Equal nonzero weights now distinguish different tangent characters, while a single coordinate supplies two dimensions at weight zero.

\begin{proof}[Proof of Theorem~\ref{thm:one-slope-all-indices}]
Choose exact tame eigenparameters with tangent scalars $\alpha_1,\alpha_2\in\mu_m$. Choose $rD\equiv1\pmod m$ and extend the tame generator by $\tau(\varpi)=\zeta^r\varpi$. Both scaling exponents are $pe$, so the scaled tangent eigenvalues are
\[
 \lambda_i=\alpha_i\zeta^{-rpe}.
\]
Suppose $\alpha_1\ne\alpha_2$. Their $j$th Frobenius powers have distinct tame eigenvalues $\lambda_1^{p^j},\lambda_2^{p^j}$, since Frobenius on $k$ is injective. Put $d=2n-2-v$. The hypothesis $n\ge v+2$ gives $d\ge n$. For $j=0,\ldots,d-1$, each of the $d$ distinct common nonzero weights therefore contributes two independent vectors. Weight zero contains a two-dimensional stable span by Lemma~\ref{os-lem:stable-spans}. Equation~\eqref{os-eq:primitive-dimension} would imply
\[
 2n\ge2d+2\ge2n+2,
\]
a contradiction. Thus $\alpha_1=\alpha_2$, which gives equality of their unordered Frobenius orbits and hence of the special characters. No numerical upper bound on $e$ was used.
\end{proof}

\subsection{Matching canonical lines}
For two two-slope models, different kernels of the leading coordinate maps provide the independence that the polygon comparison supplied above.

\begin{lemma}[A necessary condition for unmatched canonical lines]\label{os-lem:canonical-matching}
Let $n\ge1$ and let $f:G[p^n]\simeq H[p^n]$ compare two two-slope models of parameters $z_1,z_2$, with canonical $p$-lines $C_1,C_2$. Put $s_i=v_p(z_i)$. If $f(C_1)\ne C_2$, then
\[
 \min(s_1,2n-1)+\min(s_2,2n-1)\ge2n.
\]
In particular, $s_1,s_2\ge1$ and $s_1+s_2\ge2n$. Thus $s_1+s_2<2n$ forces the given isomorphism to identify the canonical lines. There is no bound on $e$ and no hypothesis that the special characters differ.
\end{lemma}
\begin{proof}
Use the scaled graph with exponents $z_1,z_2$ and $D=p^{2n-1}$. A geometric $p^n$-point $P$ has first scaled coordinate zero on specialization precisely when $p^{n-1}P\in C_1$. Indeed the minimum value $z_1/D$ is attained exactly by successive low divisions of the noncanonical first-level points. Every division of such a point remains on the low branch. There are $(p^2-p)p^{2n-2}$ such points, and the remaining points have strictly larger value, as in Lemma~\ref{os-lem:minima}. The corresponding assertion for the second coordinate is $p^{n-1}f(P)\in C_2$.

After extension to a splitting valuation ring, the scaled coordinate specialization maps are additive. Each image is a nonzero $\F_p$-line by Lemma~\ref{os-lem:stable-spans}. Choose a nonzero basis in each line. The resulting maps to $\F_p$ have kernels
\[
 \{P:p^{n-1}P\in C_1\},\qquad
 \{P:p^{n-1}P\in f^{-1}(C_2)\}.
\]
Multiplication by $p^{n-1}$ maps the common torsion group onto its two-dimensional $\F_p$-space of $p$-points. If the canonical lines do not match, these kernels are distinct. The two maps are therefore independent, and their joint image is all of $\F_p^2$.

Consequently no Frobenius power of $X$ is proportional to a Frobenius power of $Y$ in $B_0$. Evaluation on this full product would identify two independent coordinate functions: a $p$th power multiplies each chosen coordinate basis by a nonzero scalar and fixes its $\F_p$ argument. This also proves that sufficiently high powers of $X$ and $Y$, whose weights are both zero, are independent.

The two chains have respectively
\[
 d_i=\max(0,2n-1-s_i)
\]
distinct nonzero weights. Every overlapping weight contributes two independent vectors, every other weight contributes one, and weight zero contributes two. Equation~\eqref{os-eq:primitive-dimension} gives
\[
 2n\ge d_1+d_2+2
   =4n-\min(s_1,2n-1)-\min(s_2,2n-1).
\]
This is the asserted inequality. If either $s_i$ were zero, its left-hand truncated sum would be at most $2n-1$; hence both are positive. Removing the truncations gives the remaining assertions.
\end{proof}

\subsection{The full filtration class in the remaining range}
The cyclotomic determinant fixes dimensions before the mixed and polygon obstructions are applied.

\begin{proof}[Proof of Corollary~\ref{cor:all-index-two-slope-reduction}]
The assumptions and $e\ge(p-1)p^v$ give
\[
 p-1\le p^{2n-3-2v},\qquad v\le n-2.
\]
A primitive $p^n$th root satisfies
\[
 v_K(\zeta_{p^n}-1)=\frac{e}{(p-1)p^{n-1}}.
\]
This follows from the Eisenstein polynomial $\Phi_{p^n}(1+T)$ over a Cohen coefficient ring: its constant term is $p$ and its reduction is $T^{(p-1)p^{n-1}}$. The displayed value is not an integer, since its $p$-adic valuation is $v-n+1<0$. Thus the completed maximal unramified base does not contain $\mu_{p^n}$. Because it contains $\mu_p$, the cyclotomic image modulo $p^n$ is a nontrivial $p$-group, already attained by wild inertia.

The determinants have the form $\chi_p^d\epsilon$, with $0\le d\le2$ and $\epsilon$ of order dividing $p-1$. Equality of truncated determinants on wild inertia forces $p\mid d-d'$, hence $d=d'$. Reduction modulo $p$ then identifies the finite factors. Dimension zero is recovered by the residual Brauer character, and Cartier duality treats dimension two. In dimension one, uniqueness descends an ordinary connected--\'etale filtration, leaving the alternatives of two potentially good height-one graded pieces or a whole potentially supersingular group. Two filtered groups compare residually. A mixed unequal-character pair is excluded by Theorem~\ref{thm:mixed-all-indices}.

For two whole supersingular groups, a comparison involving different polygons contradicts Theorem~\ref{thm:mixed-polygon-obstruction}, because its necessary bound is larger by a factor of $p$. Two one-slope groups have equal characters by Theorem~\ref{thm:one-slope-all-indices}, since $n\ge v+2$. Only two two-slope models remain, as asserted.
\end{proof}

\subsection{Two Frobenius chains at arbitrary ramification}

We first isolate the dimension calculation that applies without a
prime-to-$p$ hypothesis on the absolute ramification index.  Retain the
common tame good field, the exact tame eigenparameters, and the scaled
actual graph constructed above.  Thus
\[
 m=p^2-1,\qquad D=p^{2n-1},\qquad
 \tau(x_i)=\alpha_i x_i,
 \qquad \rho D\equiv1\pmod m,
\]
and the scaled coordinates have tame eigenvalues
\[
 \lambda_i=\alpha_i\zeta^{-\rho z_i}.
\]
We use the same notation for prime-to-$p$ roots of unity and their
reductions.  The inverse $\rho$ is unrelated to a remaining torsion
level.

\begin{lemma}\label{os-lem:two-chains}
Suppose that two whole potentially supersingular groups have two-slope
parameters $z_1,z_2$, unequal special characters, and isomorphic actual
$p^n$-torsion.  Put $M=2n-1$, and write, after relabeling,
\[
 z_1=p^s a,\qquad z_2=p^t b,
 \qquad p\nmid ab,\qquad 0\le s\le t<M.
\]
Then $\lambda_1,\lambda_2\in\mathbf F_p^{\times}$.  If
$\lambda_1\ne\lambda_2$, then $s+t\ge2n$.  If
$\lambda_1=\lambda_2$, then $a\ne b$ and
\begin{equation}\label{os-eq:two-chain-divisibility}
 (p-1)p^{\max\{0,\,2n-1-s-t\}}\mid a-b.
\end{equation}
These conclusions impose no condition that $p\nmid e$.
\end{lemma}

\begin{proof}
The rescaled graph is a finite free order of rank $p^{2n}$, with
primitive special coordinates $X,Y$.  Its primitive space $V$ satisfies
\eqref{os-eq:primitive-dimension}, so $\dim V=2n$.  The diagonalizable
group scheme $\mu_D$ gives the weights
\[
 \operatorname{wt}(X^{p^j})=-z_1p^j\pmod D,
 \qquad
 \operatorname{wt}(Y^{p^j})=-z_2p^j\pmod D.
\]
This is a group-scheme grading, including in characteristic $p$; it
does not assume that $\mu_D$ has $D$ distinct geometric points.
Every displayed Frobenius iterate is nonzero, as in
Lemma~\ref{os-lem:stable-spans}: a vanishing iterate would make a
positive power of a rescaled coordinate divisible by the uniformizer,
contradicting evaluation at a point where that coordinate is a unit.

We give the tame leading-coefficient calculation explicitly.  By
Lemma~\ref{os-lem:minima}, a two-slope coordinate has exactly
$(p-1)D$ points of minimum value $z_i/D$, and exactly $D$ points,
including zero, of larger value.  In the extension with uniformizer
$u$, $u^D=\pi_L$, consider the monogenic order generated by
$X_i=u^{-z_i}x_i$.  It is finite free of rank $p^{2n}$ and has a
monic defining polynomial $P_i$ of that degree.  Its rescaled formal
addition has nonlinear coefficients divisible by $u$, so the
reduction of $X_i$ is primitive.  Hence $\overline P_i$ is additive.
Indeed, the mixed terms of
$\overline P_i(U+V)-\overline P_i(U)-\overline P_i(V)$ have degree
less than $\deg P_i$ in each variable; their vanishing modulo
$(\overline P_i(U),\overline P_i(V))$ forces their coefficients to
vanish.  The constant coefficient is zero because the identity is a
root.

Exactly $D$ rescaled roots have positive valuation, while the other
$(p-1)D$ roots are units.  Thus zero has multiplicity exactly $D$ in
$\overline P_i$.  Additivity, and the fact that $D$ and $pD$ are
consecutive powers of $p$, give
\[
 \overline P_i(T)=T^{pD}+b_iT^D,\qquad b_i\ne0.
\]
Tame inertia acts trivially on the residue field and sends $X_i$ to
$\lambda_iX_i$.  Comparing the two nonzero coefficients therefore
gives $\lambda_i^{(p-1)D}=1$.  Injectivity of Frobenius in the residue
field implies $\lambda_i^{p-1}=1$, proving
$\lambda_i\in\mathbf F_p^{\times}$.

The first chain has $M-s$ distinct nonzero weights, followed by
nonzero vectors of weight zero.  Its span consequently has dimension
at least $M-s+1$; the other chain similarly has dimension at least
$M-t+1$.  Since each $\lambda_i$ lies in $\mathbf F_p$, its entire
Frobenius chain belongs to the $\lambda_i$ tame eigenspace.  If the
eigenvalues differ, these spans are disjoint, and
\[
 (M-s+1)+(M-t+1)\le2n.
\]
This is $s+t\ge2n$.

Suppose henceforth that the leading eigenvalues agree.  If $a=b$,
then
\[
 \alpha_2
 =\lambda_1\zeta^{\rho p^t a}
 =\bigl(\lambda_1\zeta^{\rho p^s a}\bigr)^{p^{t-s}}
 =\alpha_1^{p^{t-s}}.
\]
The special Frobenius orbits would agree, contradicting the assumed
inequality of characters.  Hence $a\ne b$; put
$\ell=v_p(a-b)$.

A nonzero weight of the two chains can coincide only at equal
$p$-adic order $w$, where $t\le w<M$.  At that order the equality
is equivalent to
\[
 a\equiv b\pmod {p^{M-w}}.
\]
Thus at most $\min\{\ell,M-t\}$ nonzero weights are shared.  The
number of distinct represented nonzero weights, together with at
least one weight-zero vector, gives
\[
 (M-s)+(M-t)-\min\{\ell,M-t\}+1\le\dim V=M+1.
\]
It follows that $\ell\ge M-s-t$ whenever the latter integer is
positive.

It remains to retain the factor $p-1$.  The actual residual
isomorphism identifies determinant characters.  Since $\mu_p\subset
K$, the residual cyclotomic factor is trivial; reduction is injective
on prime-to-$p$ roots of unity.  The two finite determinant factors
are therefore equal.  For the special tangent characters this says
\[
 \alpha_1^{p+1}=\alpha_2^{p+1}.
\]
Equality of $\lambda_1,\lambda_2$ now implies
$\zeta^{\rho(p+1)(z_1-z_2)}=1$.  Since $\rho$ is invertible modulo
$m=(p-1)(p+1)$, we obtain $(p-1)\mid z_1-z_2$.  Reducing
$p^s a-p^t b$ modulo $p-1$ yields $(p-1)\mid a-b$.
Combining this coprime divisibility with the preceding power of $p$
proves \eqref{os-eq:two-chain-divisibility}.
\end{proof}


\subsection{The uniform bound at every index}
The two-chain divisibility forces a contradiction in the larger uniform range, without using a bound on the ramification of a chosen point field.

\begin{proof}[Proof of Theorem~\ref{thm:all-index-upper}]
Let $e\le p^{n-1}$ and put $v=v_p(e)$. Since $e\ge(p-1)p^v>p^v$, one has $v\le n-2$. Hence
\[
 n\ge v+2,\qquad e\le p^{n-1}\le p^{2n-3-v}.
\]
Corollary~\ref{cor:all-index-two-slope-reduction} treats the determinant, dimension, filtered, mixed and one-slope cases. An unequal-character comparison would therefore consist of two whole supersingular two-slope models. Their parameters satisfy
\[
 0<z_i<pe\le p^n.
\]
After relabeling write $z_1=p^s a$, $z_2=p^t b$, with $p\nmid ab$ and $0\le s\le t\le n-1<2n-1$. Lemma~\ref{os-lem:two-chains} applies at this original torsion level. Different leading eigenvalues would force $s+t\ge2n$, contrary to $s+t\le2n-2$.

For equal leading eigenvalues, the lemma gives $a\ne b$ and
\[
 (p-1)p^{2n-1-s-t}\mid a-b.
\]
The exponent is at least $n-s$, whereas positivity and the strict upper bounds on $z_i$ give
\[
 0<|a-b|<p^{n-s}.
\]
A nonzero integer of that size cannot be divisible by $(p-1)p^{n-s}$. This contradiction excludes the last case and proves the uniform recovery assertion for every $n\ge2$.

For its index formulation put $N=\lceil\log_p e\rceil$. Since $e\ge p-1>1$, one has $N\ge1$. Apply the assertion with $n=N+1$ to obtain $\nuMin(p,e,2)\le N$. The graph range in Proposition~\ref{prop:graph-uniform-full} follows because $(p-1)p^{n-2}<p^{n-1}$, and its direct proof gives an additional route in that subrange.
\end{proof}

\subsection{High Hasse valuation in the proposed range}
A sharper proposed recovery range forces at least one Hasse parameter to have two more factors of $p$ than the absolute index.

\begin{lemma}\label{os-lem:high-hasse}
Let $v=v_p(e)$, $n\ge2$, and $e\le p^{2n-3-v}$.  In an
unequal-character actual $p^n$-torsion comparison of two two-slope
whole potentially supersingular groups,
\[
 \max\{v_p(z_1),v_p(z_2)\}\ge v+2.
\]
\end{lemma}

\begin{proof}
Admissibility gives $e\ge(p-1)p^v>p^v$, so the assumed bound implies
$n\ge v+2$.  Also
\[
 0<z_i<pe\le p^{2n-2-v}<p^{2n-1}.
\]
Thus Lemma~\ref{os-lem:two-chains} applies.  Suppose, to the contrary,
that its ordered valuations satisfy $s\le t\le v+1$.
Different leading eigenvalues would require $s+t\ge2n$, whereas
$s+t\le2v+2<2n$.

For equal leading eigenvalues, write $z_1=p^sa$ and $z_2=p^tb$ as
in that lemma.  Its nonzero difference satisfies
\[
 0<|a-b|<\frac{pe}{p^s}\le p^{2n-2-v-s}.
\]
But \eqref{os-eq:two-chain-divisibility} makes that integer divisible
by $(p-1)p^{2n-1-s-t}$, and $t\le v+1$ gives
\[
 (p-1)p^{2n-1-s-t}\ge(p-1)p^{2n-2-v-s}
 >p^{2n-2-v-s}.
\]
This contradiction proves the assertion.
\end{proof}

\subsection{The central terminal comparison}

The next argument specifies exactly the comparison that remains
when the proposed sharper numerical range does not already force
recovery.  It does not retain that range after a loss of torsion
levels.

\begin{proof}[Proof of Theorem~\ref{thm:all-index-central-reduction}]
Put $v=v_p(e)$ and suppose $e\le p^{2n-3-v}$.  An
unequal-character comparison reduces by
Corollary~\ref{cor:all-index-two-slope-reduction} to two whole
potentially supersingular groups with two slopes.  By
Lemma~\ref{os-lem:high-hasse}, choose one side with parameter $z$
satisfying $v_p(z)\ge v+2$.

The minimum-value formula and Proposition~\ref{fc-thm:valuation} give
\[
 z\ge\frac{(p-1)p^{2n-1}}{p^n-1}>(p-1)p^{n-1}.
\]
Let $c$ be the least nonnegative integer with $p^cz>e$.  Such an
integer exists, and equality $p^jz=e$ never occurs: its two sides
would have different $p$-adic valuations.  Moreover
\[
 p^{n-v-2}z>(p-1)p^{2n-v-3}\ge e,
\]
so $c\le n-v-2$.

At each stage before $c$, the chosen parameter is strictly below
$e$.  The canonical quotient formula in
Proposition~\ref{rc-prop:quotients} therefore multiplies it by $p$.
Quotient the partner by the transported actual $K$-stable kernel.
The actual finite-level quotient construction gives a comparison
one level lower, and the $K$-isogenies preserve the special rational
inertia characters.  The transported kernel is not presumed to be
canonical in the partner.

After exactly $c$ such steps, the remaining level is $p^r$, where
$r=n-c\ge v+2$.  The chosen parameter is $z_c=p^cz$.  If $c=0$,
its original two-slope range gives $z_c<pe$; if $c>0$, the preceding
parameter was strictly below $e$, again giving $z_c<pe$.  Hence
\[
 e<z_c<pe.
\]
Finally $v_p(z_c)\ge v+2+c$, whereas
$v_p((p+1)e)=v$.  Therefore
\[
 v_p\bigl((p+1)e-z_c\bigr)=v.
\]
The partner remains whole potentially supersingular under these
isogenies, but its terminal polygon is unrestricted.  The original
inequality involving $n$ has not been asserted with the smaller
integer $r$.  These are exactly the stated terminal conclusions.
\end{proof}

\subsection{The full-class index and the supersingular index}

We now separate the mixed obstruction from the remaining whole
supersingular comparison.  Define $\nu_{\mathrm{SS}}(p,e)$ to be the
least integer $a\ge0$ for which actual $p^{a+1}$-torsion determines
the special character uniformly among two whole potentially
supersingular groups over all standing fields of absolute index $e$.
Its finiteness follows from the uniform upper theorem.  Both this
index and $\nu_{\min}(p,e,2)$ use universal upper quantifiers; an
existential lower example at the specified index is therefore enough
to disprove recovery at a given level.

\begin{proof}[Proof of Theorem~\ref{thm:full-supersingular-maximum}]
Write $v=v_p(e)$ and
\[
 R=\left\lceil\frac{1+\log_p e+v}{2}\right\rceil,
 \qquad T=\max\{\nu_{\mathrm{SS}}(p,e),R\}.
\]
Since $e\ge(p-1)p^v$, one has $R\ge v+1$.  For $n=T+1$ it
follows that $n\ge v+2$ and
\[
 e\le p^{2R-1-v}\le p^{2T-1-v}=p^{2n-3-v}.
\]
An unequal-character full-class comparison at level $p^n$ would
therefore reduce, by
Corollary~\ref{cor:all-index-two-slope-reduction}, to a comparison
of two whole potentially supersingular groups.  This contradicts
$T\ge\nu_{\mathrm{SS}}(p,e)$, by restriction of torsion if necessary.
Consequently $\nu_{\min}(p,e,2)\le T$.

Conversely, inclusion of the supersingular comparison class gives
$\nu_{\min}(p,e,2)\ge\nu_{\mathrm{SS}}(p,e)$.  If $R\ge2$, the
strict inequality defining the ceiling gives
\[
 2R-2<1+\log_p e+v,
 \qquad e>p^{2R-3-v}.
\]
The actual mixed existence assertion of
Theorem~\ref{thm:mixed-all-indices}, applied at level $p^R$, supplies
unequal-character groups with isomorphic actual $p^R$-torsion over
some standing field of that exact absolute index.  Hence
$\nu_{\min}(p,e,2)\ge R$.

If $R=1$, the already proved inequality $R\ge v+1$ gives $v=0$,
and the defining bound gives $e\le p$.  A positive multiple of
$p-1$ in this range must be $e=p-1$ because $p\ge5$.
Theorem~\ref{thm:minimal-index} supplies the required lower bound
$\nu_{\min}(p,p-1,2)=1$.  Thus the same lower estimate holds in
this sole case not covered by the mixed theorem's $n\ge2$ range.
Combining both lower bounds with the upper bound proves
\[
 \nu_{\min}(p,e,2)
 =\max\left\{\nu_{\mathrm{SS}}(p,e),
 \left\lceil\frac{1+\log_p e+v_p(e)}2\right\rceil\right\}.
\]
\end{proof}

\subsection{Extending the exact wild bands}\label{sec:extended-wild-bands}

We now enlarge the coefficient interval by matching the canonical
residual characters before taking one actual quotient.  The argument
first treats coefficients prime to $p$; the different values at
multiples of $p$ then follow from the smaller coefficient bands.

\begin{proof}[Completion of Theorem~\ref{thm:wild-bands-exact}]
Assume first that
\[
 2\le k\le\frac{p^2-1}{2},\qquad p\nmid k,
 \qquad e=k(p-1)p^r,
 \qquad n=r+3,\qquad N=n-1=r+2.
\]
In particular $v_p(e)=r$.  Since $k<p^2$, we have
$e<p^{r+3}=p^{2n-3-r}$.  An unequal-character actual $p^n$-torsion
comparison in the full filtration class therefore reduces, by
Corollary~\ref{cor:all-index-two-slope-reduction}, to two whole
potentially supersingular groups with two slopes.
Lemma~\ref{os-lem:minima} and Proposition~\ref{fc-thm:valuation} give
\begin{equation}\label{os-eq:half-square-range}
 (p-1)p^{n-1}<z_i<pe=k(p-1)p^{n-2}.
\end{equation}
Both models are already central, because the lower bound divided by
$e$ is $p^2/k>1$.  The upper bound is less than $p^{n+1}$, so
$v_p(z_i)\le n$.

\smallskip\noindent
\emph{A parameter with valuation $n$.}
Suppose both valuations were at most $n-1$, ordered $s\le t$.
Different leading tame eigenvalues would contradict the inequality
$s+t\ge2n$ in Lemma~\ref{os-lem:two-chains}.  For equal eigenvalues,
write $z_1=p^sa$ and $z_2=p^tb$, with $p\nmid ab$.  Unequal special
characters make $a-b$ nonzero, and the same lemma makes it divisible
by at least $(p-1)p^{n-s}$.  But
\[
 |a-b|<k(p-1)p^{n-2-s}<(p-1)p^{n-s},
\]
since $k<p^2$.  Thus one valuation is exactly $n$.  Relabel that
parameter as
\begin{equation}\label{os-eq:half-square-high}
 z_1=bp^n,
 \qquad 1\le b<\frac{k(p-1)}{p^2}<\frac{p-1}{2}.
\end{equation}
In particular $b$ is prime to $p$.  Every parameter of valuation $n$
has its coefficient in this same strict interval.

\smallskip\noindent
\emph{Matching the canonical lines.}
Let $m=p^2-1$, let $\zeta$ have order $m$, and write
$\alpha_i=\zeta^{a_i}$ for the special tangent scalars.  Put
$\delta=\zeta^{p+1}$, of order $p-1$.  The tame inverse of
$p^{2n-1}$ modulo $m$ is $p$, so the leading scalars of
Lemma~\ref{os-lem:two-chains} are
\[
 \lambda_i=\alpha_i\zeta^{-pz_i}\in\mathbf F_p^{\times}.
\]
The canonical residual character $\kappa_i$ and its quotient
character $\eta_i$ satisfy
\begin{equation}\label{os-eq:half-square-residual}
 \kappa_i=\zeta^{(p+1)e-z_i-a_i},
 \qquad
 \eta_i=\zeta^{pz_i-a_i}=\lambda_i^{-1},
 \qquad
 \frac{\kappa_i}{\eta_i}=\delta^{-z_i}.
\end{equation}
For the first formula, a canonical point has high value
$(p+1)e-z_i$, and its leading field multiplier is divided by
$\alpha_i$ under descent, as in \eqref{rc-eq:kappa}.
For the second, the minimum-point description in
Lemma~\ref{os-lem:minima} identifies the leading noncanonical
quotient multiplier with $\zeta^{pz_i}$, again divided by
$\alpha_i$.  Equivalently this is the leading-coordinate computation
of \eqref{pg-eq:residual-characters}.  These leading-term calculations
use no prime-to-$p$ condition on $e$.  The last equality uses
$(p-1)\mid e$.

Suppose $\lambda_1\ne\lambda_2$.  The valuation-sum bound in
Lemma~\ref{os-lem:two-chains} forces $v_p(z_2)=n$, so
$z_2=b'p^n$ with $0<b'<(p-1)/2$.  The two quotient characters
$\eta_i$ differ.  Equality of the unordered residual composition
characters must therefore interchange them:
\[
 \eta_2=\kappa_1,\qquad \kappa_2=\eta_1.
\]
Equation~\eqref{os-eq:half-square-residual} then gives
$(p-1)\mid z_1+z_2$, hence $(p-1)\mid b+b'$.  This contradicts
$0<b+b'<p-1$.  The leading eigenvalues consequently agree.

If the second valuation were also $n$, the two-chain congruence
would give $(p-1)\mid b-b'$.  The strict interval in
\eqref{os-eq:half-square-high} would force $b=b'$, and equal
parameters and leading scalars would give equal tangent scalars,
contrary to unequal special characters.  Thus, writing
$s=v_p(z_2)$, we have $s\le n-1$.  The two-chain congruence gives
an integer $j$ such that
\begin{equation}\label{os-eq:half-square-partner}
 z_2=bp^s+j(p-1)p^{n-1}.
\end{equation}
The lower bound in \eqref{os-eq:half-square-range} implies $j\ge1$.
No upper bound on $j$ will be needed.

Equal leading scalars give equal quotient characters $\eta_i$.
The residual determinant then gives equal canonical characters
$\kappa_i$.  Their ratio in the first model is $\delta^{-b}\ne1$,
by \eqref{os-eq:half-square-high}.  Thus both representations have
distinct canonical and quotient characters.  A two-dimensional
representation with these composition characters has at most one
stable line of character $\kappa_i$: two distinct such lines would
span the space and make both composition factors equal to
$\kappa_i$.  The actual residual isomorphism therefore identifies
the canonical lines.

\smallskip\noindent
\emph{The quotient parameters.}
Take these two matched canonical quotients.
Lemma~\ref{lem:actual-quotients} gives an actual $p^N$ comparison,
and the $K$-isogenies preserve the unequal special characters.
Both original models were proved central.  Hence
Proposition~\ref{rc-prop:quotients} gives two central, two-slope
quotient parameters
\begin{equation}\label{os-eq:half-square-quotients}
 \begin{aligned}
 A&=p^rU, & U&=k(p^2-1)-bp^3,\\
 B&=p^rV-bp^s, & V&=k(p^2-1)-j(p-1)p^2.
 \end{aligned}
\end{equation}
In particular $e<A,B<pe$.  Since $U\equiv V\equiv-k\pmod p$,
one has $v_p(A)=r$.  If $s<r$, then $v_p(B)=s$; if $s>r$, then
$v_p(B)=r$.  If $s=r$ and $p\nmid k+b$, again $v_p(B)=r$.

In each of these cases the ordered quotient valuations $u\le t$
satisfy $u,t\le r=N-2$.  They are below $2N-1$ and their sum is
less than $2N$, so Lemma~\ref{os-lem:two-chains} forces equal
leading eigenvalues.  The nonzero difference of their unit parts
is then divisible by
\[
 (p-1)p^{2N-1-u-t}\ge(p-1)p^{r+3-u}.
\]
Its absolute value is instead less than
\[
 \frac{pe}{p^u}=k(p-1)p^{r+1-u},
\]
which is strictly smaller because $k<p^2$.  This contradiction
excludes every such case.  Only the explicit valuation hypotheses
of the two-chain lemma have been used at level $N$; the original
numerical candidate range has not been carried to that lower level.

\smallskip\noindent
\emph{The cancellation case.}
It remains to treat $s=r$ and $p\mid k+b$.  Write
\[
 k=pq+c,\qquad 1\le c\le p-1.
\]
The coefficient bounds give
\[
 q\le\frac{p-1}{2},\qquad b\le\frac{p-3}{2}.
\]
Since $0<c+b<2p$, the divisibility forces $c+b=p$.  Thus
$k+b=p(q+1)$, with $1\le q+1<p$, and
\eqref{os-eq:half-square-quotients} becomes
\[
 B=p^{r+1}W,
 \qquad W=p\bigl[k-j(p-1)\bigr]-(q+1).
\]
Consequently $v_p(B)=r+1$ exactly.  The quotient valuations $r,r+1$
are below $2N-1$ and have sum less than $2N$.  The two-chain lemma
again forces equal leading eigenvalues and requires $p^2$ to divide
$U-W$.  But
\[
 U-W\equiv-k+q+1\equiv b+q+1\pmod p,
 \qquad 1\le b+q+1\le p-1.
\]
This is the final contradiction.  The uniform upper bound
$\nuMin(p,e,2)\le r+2$ has been proved for the prime-to-$p$
coefficients in the stated interval.

\smallskip\noindent
\emph{Lower examples and coefficients divisible by $p$.}
For every such prime-to-$p$ coefficient, the strict mixed threshold
at level $p^{r+2}$ is satisfied:
\[
 e=k(p-1)p^r>p^{r+1}=p^{2(r+2)-3-r},
\]
because $k\ge2$ and $p\ge5$.  Theorem~\ref{thm:mixed-all-indices}
therefore supplies actual unequal-character groups with isomorphic
$p^{r+2}$-torsion over some standing field of the exact index $e$,
which may have residue field $\overline{\mathbf F}_p$.  It follows
that $\nuMin(p,e,2)\ge r+2$, proving equality.

Finally suppose $p\mid k$ in the same coefficient interval.  Write
$k=jp$.  Since $k\le(p^2-1)/2$, one has
$1\le j\le(p-1)/2$.  If $j=1$, the already established coefficient
$k=p$ case in Section~\ref{sec:all-wild-indices} gives
\[
 \nuMin\bigl(p,p^{r+1}(p-1),2\bigr)=r+2.
\]
If $j\ge2$, apply the coefficient-$j$ case proved in that section,
with exponent $r+1$, to obtain
\[
 \nuMin\bigl(p,jp^{r+1}(p-1),2\bigr)=r+3.
\]
These earlier small-coefficient proofs include both their uniform
upper statements and their actual existential lower constructions.
They do not use the prime-to-$p$ valuation calculation of the present
argument.  This proves every case of the theorem, with the stated
universal upper and existential lower quantifiers.
\end{proof}
